\chapter{Aljabar Linear}\label{ch:b2:linalg}

Aljabar linear pada jilid Tahun ke-1 bekerja atas $\R$ atau $\C$ dalam
dimensi hingga, dan menerima \omterm{def:b2:linalg:det}{determinan} umum tanpa bukti. Bab ini
menaikkan ketiga pembatasan itu: teorinya dinyatakan atas lapangan
$K$ sembarang, jalinan antara sebuah ruang dan \emph{dual}-nya
dikembangkan secara sistematis (\omterm{def:b2:linalg:dual}{basis dual}, \omterm{def:b2:linalg:annihilator}{anihilator}, \omterm{def:b2:linalg:transpose}{transpos}), dan
\omterm{def:b2:linalg:det}{determinan} akhirnya \emph{\omterm{def:b2:structures:generated}{dibangun}} dari bentuk multilinear \omterm{def:b2:linalg:alternating}{alternating}
serta tanda permutasi pada
\cref{ch:b2:structures} --- sehingga melunasi setiap utang Tahun ke-1.

Di sepanjang bab ini, $K$ adalah lapangan ($\Q$, $\R$, $\C$, atau
$\Z/p\Z$ --- teorinya tidak peduli) dan, kecuali disebutkan lain, ruangnya
berdimensi hingga atas $K$. Hasil Tahun ke-1 (basis, dimensi,
rank--nulitas, matriks) berpindah kata demi kata: buktinya tak pernah
memakai apa pun selain aksioma lapangan.

\section{Ruang dual}

\begin{definition}[Ruang dual, basis dual]\label{def:b2:linalg:dual}
\emph{Dual}\index{ruang dual} dari $E$ adalah $E^* = \mathcal{L}(E,
K)$, yaitu ruang semua bentuk linear. Jika $\mathcal{B} = (e_1, \dots, e_n)$
basis $E$, maka \emph{bentuk koordinat} $e_1^*, \dots, e_n^*$
yang ditetapkan oleh $e_i^*(e_j) = \delta_{ij}$ (Kronecker: $1$ bila $i = j$,
selain itu $0$) membentuk \emph{basis dual}\index{basis dual}
$\mathcal{B}^*$ bagi $E^*$; khususnya $\dim E^* = \dim E$, dan
\[
x = \sum_{i=1}^{n} e_i^*(x)\, e_i \quad (x \in E),
\qquad
\varphi = \sum_{i=1}^{n} \varphi(e_i)\, e_i^* \quad (\varphi \in
E^*).
\]
\end{definition}

\begin{proof}[Bukti bahwa $\mathcal{B}^*$ adalah basis]
Bebas: menerapkan kombinasi nol $\sum \lambda_i e_i^* = 0$ pada
$e_j$ memberi $\lambda_j = 0$. Membangun: untuk $\varphi \in E^*$, bentuk
$\varphi - \sum_i \varphi(e_i) e_i^*$ menolkan setiap $e_j$, jadi
ia nol (pemetaan linear yang nol pada suatu basis adalah pemetaan nol). Kedua
rumus di atas adalah perhitungan yang sama, dibaca maju.
\end{proof}

\begin{example}\label{ex:b2:linalg:dualexamples}
Pada $K_n[X]$ dengan basis $(1, X, \dots, X^n)$: \omterm{def:b2:linalg:dual}{basis dualnya} adalah $P
\mapsto \frac{P^{(k)}(0)}{k!}$ (koefisien Taylor). Basis lain
bagi dualnya: evaluasi $P \mapsto P(x_i)$ pada $n + 1$ titik
yang berbeda --- basis ``pra-dual''-nya di $K_n[X]$ persis keluarga
polinomial Lagrange $L_i$ (jilid Tahun ke-1), sebab $L_i(x_j) =
\delta_{ij}$. Interpolasi \emph{adalah} dualitas.
\end{example}

\begin{method}[Basis dual dan antedual dalam praktik]\label{met:b2:linalg:dualbasis}
Untuk menguraikan sebuah bentuk $\varphi$ pada basis $(e_i)$ di $E$:
koordinatnya adalah \emph{nilai} $\varphi(e_i)$ --- tanpa sistem yang
perlu diselesaikan. Untuk mencari basis $(u_j)$ di $E$ yang dualnya
adalah basis $(\varphi_1, \dots, \varphi_n)$ tertentu di $E^*$ (yaitu
basis \emph{antedual}): selesaikan $n$ sistem linear
\[
\varphi_i(u_j) = \delta_{ij} \qquad (1 \leq i \leq n),
\]
satu kolom $u_j$ setiap kali; dalam bahasa matriks, bila baris $M$
mendaftar koefisien $\varphi_i$ pada basis $E^*$ yang sudah diketahui,
maka kolom $M^{-1}$ adalah $u_j$. Keberadaan dan
ketunggalan antedual dibuktikan pada soal akhir pekan bab ini;
perhitungannya selalu berupa pembalikan matriks ini.
\end{method}

\begin{example}[Sebuah basis dual di $\R^2$, dihitung tuntas]\label{ex:b2:linalg:dualbasis2d}
Untuk basis $b_1 = (1, 1)$, $b_2 = (1, -1)$ di $\R^2$:
\omterm{def:b2:linalg:dual}{basis dual} $(b_1^*, b_2^*)$ wajib memenuhi $b_i^*(b_j) =
\delta_{ij}$. Dengan menulis $b_1^*(x, y) = \alpha x + \beta y$, syarat
$\alpha + \beta = 1$ dan $\alpha - \beta = 0$ memberi
\[
b_1^*(x, y) = \frac{x + y}{2},
\qquad\text{dan serupa itu}\qquad
b_2^*(x, y) = \frac{x - y}{2} .
\]
Pemeriksaan kewarasan: $b_1^*$ \emph{bukan} $e_1^* + e_2^*$ yang dinilai
begitu saja --- \omterm{def:b2:linalg:dual}{basis dual} bergantung pada seluruh basis, bukan pada
tiap vektor secara terpisah (mengganti $b_2$ dengan $(0, 1)$ mengubah
$b_1^*$ menjadi $x \mapsto x$). Dan rumus penguraiannya berjalan:
$(x, y) = \frac{x+y}2\,b_1 + \frac{x-y}2\,b_2$, yaitu penguraian
genap/ganjil sebuah pasangan --- \omterm{def:b2:linalg:dual}{basis dual} adalah pengekstrak
koordinat, dan yang satu ini mengekstrak bagian setangkup dan bagian
antisetangkup.
\end{example}

\begin{definition}[Anihilator]\label{def:b2:linalg:annihilator}
Untuk subruang $F \subseteq E$,
\emph{anihilator}\index{anihilator} adalah
\[
F^{\circ} = \{\varphi \in E^* : \varphi|_F = 0\},
\]
yaitu subruang $E^*$.
\end{definition}

\begin{theorem}[Dimensi anihilator]\label{thm:b2:linalg:annihilator}
$\dim F^{\circ} = \dim E - \dim F$. Lebih lanjut $F \mapsto F^\circ$
membalik pemuatan, dan $F$ dapat dipulihkan dari \omterm{def:b2:linalg:annihilator}{anihilatornya}:
\[
F = \{x \in E : \forall\varphi \in F^\circ,\ \varphi(x) = 0\}.
\]
Akibatnya setiap subruang berdimensi $p$ di dalam dimensi $n$ adalah
himpunan penyelesaian $n - p$ persamaan linear yang bebas --- dan
sebaliknya.
\end{theorem}

\begin{proof}
Pilih basis $(e_1, \dots, e_p)$ untuk $F$, lalu lengkapi menjadi basis
$E$. Sebuah bentuk $\varphi = \sum \varphi(e_i) e_i^*$ menganihilasi $F$
bila dan hanya bila $p$ koefisien pertamanya nol: jadi $F^\circ =
\operatorname{Vect}(e_{p+1}^*, \dots, e_n^*)$, berdimensi $n - p$.
Pembalikan pemuatan langsung terlihat. Untuk pemulihannya: ruas kanan
memuat $F$; sebaliknya, jika $x \notin F$, lengkapilah basis
$F$ dengan $x$ dan vektor lain; bentuk koordinat $x$ pada basis ini
menganihilasi $F$ tetapi tidak menganihilasi $x$. Pembacaan ``persamaan''
mengambil basis $(\varphi_1, \dots, \varphi_{n-p})$ untuk $F^\circ$: maka $F =
\bigcap \ker\varphi_j$, yaitu irisan $n - p$ hiperbidang yang bebas.
\end{proof}

\begin{example}[Sebuah anihilator, dua arah]\label{ex:b2:linalg:annihilatorwork}
Misalkan $F = \operatorname{Vect}\bigl((1, 2, 1),\ (1, 0, -1)\bigr)
\subseteq \R^3$. Sebuah bentuk $\varphi = a\,e_1^* + b\,e_2^* + c\,e_3^*$
menganihilasi $F$ bila dan hanya bila
\[
a + 2b + c = 0
\qquad\text{dan}\qquad
a - c = 0 ,
\]
yakni $c = a$ dan $b = -a$, jadi $F^\circ = \R\,(e_1^* - e_2^* +
e_3^*)$, berdimensi $3 - 2 = 1$ sebagaimana dituntut
\cref{thm:b2:linalg:annihilator}. Dibaca terbalik:
$F = \{(x, y, z) : x - y + z = 0\}$ --- bidangnya dipulihkan sebagai
kernel satu-satunya bentuk yang merentang $F^\circ$. Beranjak dari
keluarga perentang ke persamaan \emph{adalah} menghitung \omterm{def:b2:linalg:annihilator}{anihilator};
beranjak dari persamaan ke parametrisasi adalah menghitung
\omterm{def:b2:linalg:annihilator}{pra-anihilator}. (Periksa: kedua vektor perentangnya memenuhi $x - y +
z = 0$.)
\end{example}

\begin{definition}[Pemetaan transpos]\label{def:b2:linalg:transpose}
Untuk $u \in \mathcal{L}(E, F)$, \emph{transpos}\index{pemetaan
transpos} $u^{\mathsf T} \in \mathcal{L}(F^*, E^*)$ adalah
\[
u^{\mathsf T}(\psi) = \psi \circ u .
\]
Transpos memenuhi $(v \circ u)^{\mathsf T} = u^{\mathsf T} \circ v^{\mathsf
T}$, dan pada \omterm{def:b2:linalg:dual}{basis dual}, matriks $u^{\mathsf T}$ adalah
matriks transpos $u$ --- yang akhirnya \emph{menjelaskan}
transpos pada Tahun ke-1.
\end{definition}

\begin{example}[Transpos, unsur demi unsur]\label{ex:b2:linalg:transposework}
Misalkan $u \colon \R^2 \to \R^3$ bermatriks $A =
\begin{psmallmatrix} 1 & 2\\ 0 & 1\\ 3 & 0\end{psmallmatrix}$ pada
basis kanonik. Untuk $\psi = b_1f_1^* + b_2f_2^* + b_3f_3^*
\in (\R^3)^*$, hitung $u^{\mathsf T}(\psi) = \psi \circ u$ pada
basis $\R^2$:
\[
(\psi \circ u)(e_1) = \psi(1, 0, 3) = b_1 + 3b_3,
\qquad
(\psi \circ u)(e_2) = \psi(2, 1, 0) = 2b_1 + b_2 .
\]
Jadi $u^{\mathsf T}(\psi) = (b_1 + 3b_3)\,e_1^* + (2b_1 +
b_2)\,e_2^*$, dan pada \omterm{def:b2:linalg:dual}{basis dual} matriks $u^{\mathsf T}$
adalah
\[
\begin{pmatrix} 1 & 0 & 3\\ 2 & 1 & 0\end{pmatrix}
= A^{\mathsf T} :
\]
\omterm{def:b2:linalg:transpose}{transpos} abstrak \emph{adalah} matriks yang dibalik, tanpa
perhitungan yang tersisa untuk dipercaya begitu saja. Perhatikan
mekanismenya: \emph{kolom} ke-$j$ pada $A$ menjadi \emph{baris} ke-$j$
pada matriks yang baru, sebab $\psi \circ u$ membaca keluaran $u$ lewat
koefisien $\psi$.
\end{example}

\begin{proposition}\label{prop:b2:linalg:transposerank}
$\ker u^{\mathsf T} = (\operatorname{im} u)^{\circ}$ dan
$\operatorname{im} u^{\mathsf T} = (\ker u)^{\circ}$. Akibatnya
$\operatorname{rk}(u^{\mathsf T}) = \operatorname{rk}(u)$: rank baris
sama dengan rank kolom, dibuktikan secara struktural.
\end{proposition}

\begin{proof}
$\psi \in \ker u^{\mathsf T} \iff \psi \circ u = 0 \iff \psi$ menolkan
$\operatorname{im} u$: itulah kesamaan pertama. Untuk yang kedua:
$u^{\mathsf T}(\psi) = \psi \circ u$ selalu menolkan $\ker u$, jadi
$\operatorname{im} u^{\mathsf T} \subseteq (\ker u)^\circ$;
dimensinya berimpit menurut rank--nulitas dan
\cref{thm:b2:linalg:annihilator}:
\[
\operatorname{rk} u^{\mathsf T} = \dim F^* - \dim\ker u^{\mathsf T}
= \dim F - \bigl(\dim F - \operatorname{rk} u\bigr)
= \operatorname{rk} u
= \dim (\ker u)^{\circ} . \qedhere
\]
\end{proof}

\begin{example}[Rank dibaca dari kedua sisi]\label{ex:b2:linalg:rankbothsides}
Misalkan
\[
A = \begin{pmatrix}
1 & 2 & 0 & 1\\
0 & 1 & 1 & 1\\
1 & 3 & 1 & 2
\end{pmatrix} .
\]
\emph{Rank kolom:} baris ketiga adalah jumlah dua baris pertama, jadi
$\operatorname{rk} A \leq 2$; kolom $1$ dan $2$ bebas, sehingga
$\operatorname{rk} A = 2$. \emph{Kernel transposnya:} menyelesaikan
$A^{\mathsf T}y = 0$ memberi $y \in \R\,(1, 1, -1)$, jadi $\ker
A^{\mathsf T}$ berdimensi $1 = 3 - 2$: persis
$(\operatorname{im} A)^\circ$ di bawah penyamaan
$(\R^3)^*$ dengan vektor baris, sebagaimana ditegaskan
\cref{prop:b2:linalg:transposerank} --- satu-satunya
kaitan ``baris$_3$ = baris$_1$ + baris$_2$'' \emph{adalah}
\omterm{def:b2:linalg:annihilator}{anihilator} ruang kolomnya. Rank baris ($2$ baris yang bebas) dan
rank kolom berimpit bukan karena kebetulan, melainkan karena keduanya
sama dengan $\operatorname{rk} A = \operatorname{rk} A^{\mathsf T}$.
\end{example}

\begin{example}[Dualitas membaca sebuah aturan kuadratur]\label{ex:b2:linalg:quadratureread}
Mengapa aturan seperti aturan Simpson (\cref{exo:b2:linalg:4}) ada
dan mengapa ia tunggal? Dualitas menjawabnya sebelum perhitungan apa pun.
Pada $E = \R_2[X]$, integral $P \mapsto \int_0^1 P$ adalah satu
vektor tertentu di dalam dual $E^*$ yang berdimensi tiga; sedangkan
evaluasi di $0$, $\frac12$, $1$ membentuk \emph{basis}
$E^*$; karena itu integralnya terurai secara tunggal atas basis itu ---
dan penguraian itu \emph{adalah} aturan Simpson, lengkap dengan
koefisiennya. Pencacahan dimensi juga menakar harapan kita: pada $\R_3[X]$,
empat dimensi bentuk pada umumnya tak dapat direntang oleh tiga
evaluasi, jadi ketepatan pada polinomial kubik bukanlah hutang
dualitas; bahwa aturan Simpson itu toh mengintegralkan kubik secara tepat
adalah bonus kesetangkupan (pencoretan derajat ganjil di sekitar
$\frac12$) yang harus diperiksa dengan tangan. Aturan dengan $n + 1$
simpul adalah penguraian bentuk pengintegralan pada suatu basis
evaluasi di $\R_n[X]^*$: keberadaan dan ketunggalannya hanya berharga
satu teorema \omterm{def:b2:linalg:dual}{basis dual}; hanya derajat bonusnya yang berharga kerja.
\end{example}

\section{Bentuk multilinear alternating}

\begin{definition}\label{def:b2:linalg:alternating}
Pemetaan $f \colon E^n \to K$ disebut \emph{$n$-linear} bila ia linear
pada tiap peubahnya, dan disebut
\emph{alternating}\index{bentuk alternating} bila
ia nol setiap kali dua argumennya sama. Sifat alternating mengakibatkan
\emph{antisetangkup}: menukar dua argumen mengubah tandanya
(uraikan $f(\dots, x + y, \dots, x + y, \dots) = 0$); lebih umum lagi,
untuk $\sigma \in \mathfrak{S}_n$,
\[
f(x_{\sigma(1)}, \dots, x_{\sigma(n)}) =
\varepsilon(\sigma)\, f(x_1, \dots, x_n),
\]
dengan menguraikan $\sigma$ atas \omterm{def:b2:structures:sn}{transposisi}
(\cref{thm:b2:structures:signature}).
\end{definition}

\begin{theorem}[Teorema dasar determinan]\label{thm:b2:linalg:detspace}
Misalkan $\dim E = n$ dan $\mathcal{B} = (e_1, \dots, e_n)$ sebuah basis.
Ruang bentuk $n$-linear \omterm{def:b2:linalg:alternating}{alternating} pada $E$ berdimensi $1$:
setiap bentuk semacam itu adalah kelipatan
\[
\det{}_{\mathcal{B}}(x_1, \dots, x_n)
= \sum_{\sigma \in \mathfrak{S}_n} \varepsilon(\sigma)
\prod_{i=1}^{n} a_{\sigma(i),\,i},
\qquad
x_j = \sum_{i} a_{ij} e_i ,
\]
dan $\det_{\mathcal{B}}$ satu-satunya yang bernilai $1$ pada
$\mathcal{B}$.\index{determinan!konstruksi}
\end{theorem}

\begin{proof}
Misalkan $f$ bentuk $n$-linear \omterm{def:b2:linalg:alternating}{alternating}. Dengan menguraikan tiap
argumen pada $\mathcal{B}$ lewat kemultilinearan,
\[
f(x_1, \dots, x_n)
= \sum_{i_1, \dots, i_n} a_{i_1,1}\cdots a_{i_n,n}\,
f(e_{i_1}, \dots, e_{i_n}).
\]
Suku dengan indeks berulang bernilai nol (sifat \omterm{def:b2:linalg:alternating}{alternating}); tupel
$(i_1, \dots, i_n)$ yang bertahan adalah yang injektif, yakni $i_k =
\sigma(k)$ untuk suatu permutasi $\sigma$, dan sifat antisetangkup
menata ulang $f(e_{\sigma(1)}, \dots, e_{\sigma(n)}) = \varepsilon(\sigma)
f(e_1, \dots, e_n)$. Karenanya
\[
f = f(e_1, \dots, e_n) \cdot \det{}_{\mathcal{B}} :
\]
setiap \omterm{def:b2:linalg:alternating}{bentuk alternating} memang kelipatan itu, asalkan
$\det_{\mathcal{B}}$ sendiri (jumlah yang tertulis di atas)
\emph{memang} $n$-linear \omterm{def:b2:linalg:alternating}{alternating} dan bernilai $1$ pada $\mathcal B$.
Kemultilinearannya jelas (tiap sukunya linear pada tiap kolom).
Nilai pada $\mathcal B$: satu-satunya suku tak nol adalah $\sigma =
\mathrm{id}$. Sifat \omterm{def:b2:linalg:alternating}{alternating}: andaikan $x_j = x_k$ ($j \neq k$),
sehingga kolom koordinatnya memenuhi $a_{i j} = a_{i k}$ untuk setiap
$i$. Pasangkan tiap $\sigma$ dengan $\sigma' = \sigma\circ(j\,k)$ ---
sebuah involusi tanpa titik tetap pada $\mathfrak{S}_n$. Hasil kali
yang berpasangan itu berimpit:
\[
\prod_i a_{\sigma'(i),\,i}
= a_{\sigma(k),\,j}\; a_{\sigma(j),\,k}
\prod_{i \neq j,k} a_{\sigma(i),\,i}
= a_{\sigma(k),\,k}\; a_{\sigma(j),\,j}
\prod_{i \neq j,k} a_{\sigma(i),\,i}
= \prod_i a_{\sigma(i),\,i},
\]
dengan memakai kesamaan kolom $j$ dan kolom $k$; sementara itu
$\varepsilon(\sigma') = -\varepsilon(\sigma)$. Jadi tiap pasangan
menyumbang nol: jumlahnya lenyap.
\end{proof}

\begin{example}[Sarrus, diturunkan lalu dirobohkan]\label{ex:b2:linalg:sarrus}
Untuk $n = 3$ rumus permutasinya tepat punya $3! = 6$ suku.
Mendaftar $\mathfrak{S}_3$ menurut tandanya ---
$\mathrm{id}$, $(1\,2\,3)$, $(1\,3\,2)$ genap;
$(1\,2)$, $(1\,3)$, $(2\,3)$ ganjil --- memberi
\[
\det A = a_{11}a_{22}a_{33} + a_{21}a_{32}a_{13} +
a_{31}a_{12}a_{23}
- a_{21}a_{12}a_{33} - a_{31}a_{22}a_{13} -
a_{11}a_{32}a_{23} :
\]
persis aturan ``diagonal'' Sarrus yang diajarkan di sekolah ---
kini menjadi teorema, dengan tanda-tandanya yang misterius dikenali
sebagai tanda permutasi. Perobohannya: untuk $n = 4$ ada $24$
permutasi, dan hanya $8$ di antaranya yang terjaring oleh skema
penarikan diagonal mana pun; Sarrus tidak punya versi berderajat $4$,
dan penguraian kofaktor (\cref{thm:b2:linalg:detrules}~(4)) mengambil
alih. Mencacah sukunya sekaligus menjadi peringatan: rumus permutasi
punya $n!$ suku, jadi ia sebuah \emph{definisi}, bukan
algoritme --- penyederhanaan baris menghitung $\det$ dengan $O(n^3)$
operasi saja.
\end{example}

\begin{definition}[Determinan]\label{def:b2:linalg:det}
\emph{Determinan sebuah keluarga}\index{determinan} pada suatu basis
adalah $\det_{\mathcal{B}}(x_1, \dots, x_n)$; \emph{determinan sebuah
matriks} $A$ adalah determinan kolomnya pada basis kanonik
--- yakni rumus permutasi di atas; sedangkan \emph{determinan sebuah
endomorfisma} $u$ adalah skalar $\det u$ sedemikian sehingga
\[
\det{}_{\mathcal{B}}\bigl(u(x_1), \dots, u(x_n)\bigr)
= \det u \cdot \det{}_{\mathcal{B}}(x_1, \dots, x_n)
\quad \text{untuk setiap } x_i
\]
(ruas kirinya $n$-linear \omterm{def:b2:linalg:alternating}{alternating}, jadi merupakan kelipatan
$\det_\mathcal{B}$ menurut \cref{thm:b2:linalg:detspace}; faktornya tidak
bergantung pada $\mathcal{B}$).
\end{definition}

\begin{theorem}[Kalkulus determinan, dibuktikan]\label{thm:b2:linalg:detrules}
\begin{enumerate}
  \item $\det(uv) = \det u\,\det v$; $\;\det(AB) = \det A \det B$.
  \item $u$ punya invers $\iff \det u \neq 0$; sebuah keluarga adalah
        basis $\iff$ \omterm{def:b2:linalg:det}{determinannya} pada suatu basis tidak nol.
  \item $\det(A^{\mathsf T}) = \det A$.
  \item Penguraian kofaktor sepanjang baris atau kolom mana pun, seperti
        dinyatakan pada jilid Tahun ke-1, tetap berlaku; matriks yang
        serupa punya \omterm{def:b2:linalg:det}{determinan} yang sama.
\end{enumerate}
\end{theorem}

\begin{proof}
(1) Terapkan relasi pendefinisinya dua kali:
$\det_{\mathcal B}(uv(x_i)) = \det u \cdot \det_{\mathcal
B}(v(x_i)) = \det u \det v \cdot \det_{\mathcal B}(x_i)$.

(2) Jika $u$ punya invers, maka $\det u \det u^{-1} = \det \mathrm{id} =
1 \neq 0$. Jika tidak, peta $u(e_i)$ saling terkait; dengan menyatakan
salah satunya lewat yang lain lalu menguraikannya, $\det_{\mathcal B}(u(e_i)) = 0$
(sifat \omterm{def:b2:linalg:alternating}{alternating} menolkan arah yang berulang), jadi $\det u = 0$. Kriteria
basisnya adalah pernyataan yang sama untuk keluarga.

(3) Pada rumus permutasi, indeks ulang tiap hasil kali dengan $j =
\sigma(i)$, yakni $i = \tau(j)$ dengan $\tau = \sigma^{-1}$: faktornya
bilangan yang sama dalam urutan berbeda, sehingga
\[
\prod_{i=1}^{n} a_{\sigma(i),\,i} = \prod_{j=1}^{n}
a_{j,\,\tau(j)} ,
\]
dan $\varepsilon(\tau) = \varepsilon(\sigma)^{-1} =
\varepsilon(\sigma)$ (nilainya $\pm1$, dan $\varepsilon$ sebuah
morfisma). Menjumlahkan atas $\sigma$ sama saja dengan menjumlahkan atas
$\tau$ (pembalikan adalah bijeksi $\mathfrak{S}_n$):
\[
\det A = \sum_{\tau}\varepsilon(\tau)\prod_j a_{j,\tau(j)}
= \det(A^{\mathsf T}),
\]
sebab jumlah terakhir itu tak lain rumus permutasi yang diterapkan pada
unsur \omterm{def:b2:linalg:transpose}{transpos} $(A^{\mathsf T})_{ij} = a_{ji}$.

(4) Tetapkan kolom $j$ lalu pecah $x_j = \sum_i a_{ij} e_i$ lewat
kelinearan: $\det A = \sum_i a_{ij}\, \det(\dots, e_i, \dots)$, dan
memindahkan $e_i$ ke kedudukan terakhir ($n - i$ \omterm{def:b2:structures:sn}{transposisi} baris dan
$n - j$ \omterm{def:b2:structures:sn}{transposisi} kolom, lewat (3)) menyamakan $\det(\dots, e_i, \dots) =
(-1)^{i+j}\Delta_{ij}$ dengan minornya: persis aturan kofaktor Tahun
ke-1. Keserupaan: $\det(P^{-1}AP) = \det P^{-1}\det A \det P =
\det A$ menurut (1).
\end{proof}

\begin{example}[Penguraian kofaktor, dikerjakan]\label{ex:b2:linalg:cofactorwork}
Hitung
\[
\det\begin{pmatrix}
2 & 1 & 3\\
0 & 4 & 1\\
1 & 2 & 0
\end{pmatrix}
\]
sepanjang kolom pertama (kemalasan senilai dua nol: satu perhitungan).
Tandanya mengikuti pola papan catur $(-1)^{i+j}$:
\[
2\,\det\begin{pmatrix}4 & 1\\ 2 & 0\end{pmatrix}
- 0
+ 1\cdot\det\begin{pmatrix}1 & 3\\ 4 & 1\end{pmatrix}
= 2(0 - 2) + (1 - 12) = -15 .
\]
Periksa silang lewat Sarrus (\cref{ex:b2:linalg:sarrus}): $0 + 1 + 0
- 12 - 0 - 4 = -15$. Ini siasat, bukan doktrin: uraikan sepanjang
garis dengan nol terbanyak, dan bila tidak ada satu pun yang punya
nol, ciptakan dulu beberapa lewat operasi baris --- karena satu putaran
penyederhanaan lebih murah daripada dua lapis kofaktor.
\end{example}

\begin{example}[Sebuah determinan lewat aturannya]\label{ex:b2:linalg:onesmatrix}
Misalkan $J \in \mathcal{M}_n(K)$ matriks yang semua unsurnya satu dan
$a \in K$; kita hitung $\det(aI_n + J)$ dengan perkakas yang baru saja
dibuktikan. Setiap kolom $aI_n + J$ berjumlah dengan pola yang sama:
tambahkan semua baris ke baris pertama (\omterm{def:b2:linalg:det}{determinannya} tak berubah ---
menambahkan kelipatan satu baris ke baris lain menambahkan suku berarah
berulang, yang ditolkan sifat \omterm{def:b2:linalg:alternating}{alternating}). Baris pertama menjadi
$(a + n, a + n, \dots, a + n)$; keluarkan faktor $a + n$ lewat
kelinearan pada baris itu, lalu kurangkan kolom pertama dari setiap
kolom lain: yang tersisa segitiga dengan diagonal $(1, a, \dots, a)$. Jadi
\[
\det(aI_n + J) = (a + n)\,a^{\,n-1}.
\]
Pelajaran penutupnya: akar $a = 0$ (bermultiplisitas $n - 1$) dan
$a = -n$ mengatakan bahwa $J$ punya nilai eigen $0$ bermultiplisitas
$n - 1$ dan nilai eigen $n$ sekali --- yakni spektrum matriks $J$ yang
berank satu, satu bab lebih awal (\cref{ch:b2:reduction} akan
menjadikannya sistematis).
\end{example}

\begin{example}[Sebuah determinan lewat rumus permutasi]\label{ex:b2:linalg:permexample}
Untuk matriks dengan banyak nol, rumus itu praktis dengan sendirinya:
pada
\[
A = \begin{pmatrix}
0 & a & 0 & 0\\
0 & 0 & b & 0\\
0 & 0 & 0 & c\\
d & 0 & 0 & 0
\end{pmatrix},
\]
satu-satunya permutasi yang memungut unsur tak nol adalah \omterm{def:b2:structures:sn}{siklus}
berpanjang $4$, yaitu $\sigma = (1\,2\,3\,4)$, yang memetakan kolom $1
\to$ baris $4$, dan seterusnya; $\varepsilon(\sigma) = (-1)^3 = -1$,
jadi $\det A = -abcd$. (Periksa lewat tiga penukaran kolom sampai
tercapai matriks diagonal.)
\end{example}

\begin{example}[Vandermonde lewat rumus hasil kali]\label{ex:b2:linalg:vandermondework}
Untuk simpul $0, 1, 2$ (yang dipakai aturan kuadratur seperti pada
\cref{exo:b2:linalg:4}), \omterm{def:b2:linalg:det}{determinan} Vandermonde pada
\cref{exo:b2:linalg:11} terhitung dalam sekali pandang:
\[
\det\begin{pmatrix}
1 & 1 & 1\\
0 & 1 & 2\\
0 & 1 & 4
\end{pmatrix}
= (1 - 0)(2 - 0)(2 - 1) = 2 ,
\]
dan lewat penguraian langsung sepanjang kolom pertama: $1\cdot(4 - 2)
= 2$, jadi cocok. Ketaknolan untuk simpul yang berbeda adalah seluruh
teori interpolasi di dalam satu \omterm{def:b2:linalg:det}{determinan}: bentuk evaluasi
$P \mapsto P(a_i)$ merupakan \omterm{def:b2:linalg:dual}{basis dualnya} persis ketika
\omterm{def:b2:linalg:det}{determinan} ini tak nol, yakni selalu untuk $a_i$ yang berbeda
--- \cref{ex:b2:linalg:dualexamples} yang dikuantifikasi.
\end{example}

\section{Trace, ditinjau ulang}

\begin{proposition}\label{prop:b2:linalg:trace}
Trace $\operatorname{tr} \colon \mathcal{M}_n(K) \to K$ adalah
satu-satunya bentuk linear dengan $\operatorname{tr}(AB) =
\operatorname{tr}(BA)$ dan $\operatorname{tr}(I_n) = n$ (untuk
$\operatorname{char} K = 0$); trace sebuah endomorfisma
terdefinisi dengan baik lewat penyajian matriks mana pun, dan
\[
\operatorname{tr}(u) = \sum_{i} e_i^*\bigl(u(e_i)\bigr)
\]
pada basis mana pun --- dualitas menuliskan trace tanpa basis.
\end{proposition}

\begin{proof}
$\operatorname{tr}(AB) = \operatorname{tr}(BA)$ dan ketakbergantungannya
pada basis telah dibuktikan pada Tahun ke-1. Ketunggalan: bentuk linear
$t$ dengan $t(AB) = t(BA)$ menolkan setiap komutator $AB - BA$. Kita
klaim komutator merentang hiperbidang bertrace nol, yang berdimensi
$n^2 - 1$. Dua keluarga komutator sudah cukup. Aturan perkalian
matriks elementer adalah $E_{ab}E_{cd} = \delta_{bc}E_{ad}$. Untuk
$i \neq j$ aturan itu memberi
\[
E_{ii}E_{ij} - E_{ij}E_{ii} = E_{ij} - 0 = E_{ij}
\]
(hasil kali kedua adalah $E_{ij}E_{ii} = \delta_{ji}E_{ii} = 0$
karena $j \neq i$): jadi setiap $E_{ij}$ di luar diagonal adalah
komutator. Dan
\[
E_{ij}E_{ji} - E_{ji}E_{ij} = E_{ii} - E_{jj} .
\]
Matriks $E_{ij}$ ($i \neq j$, sebanyak $n^2 - n$) bersama
$E_{11} - E_{jj}$ ($j \geq 2$, sebanyak $n - 1$) berjumlah $n^2 - 1$
matriks bertrace nol yang bebas linear: keduanya merentang
hiperbidang $\ker\operatorname{tr}$. Jadi $t$ nol di tempat
$\operatorname{tr}$ nol, sehingga ia terfaktorkan lewatnya: $t =
c\operatorname{tr}$; lalu $t(I) = n$ memaksa $c = 1$. Adapun ungkapan
di atas: unsur diagonal ke-$i$ pada matriks $u$ tepat sama dengan
$e_i^*(u(e_i))$.
\end{proof}

\begin{remark}[Jebakan yang sering muncul]\label{rem:b2:linalg:pitfalls}
(i) \omterm{def:b2:linalg:det}{Determinan} bersifat $n$-linear pada \emph{kolomnya}, bukan
linear pada matriksnya: $\det(A + B) \neq \det A + \det B$ pada
umumnya, dan $\det(\lambda A) = \lambda^n\det A$, bukan
$\lambda\det A$. (ii) \omterm{def:b2:structures:sn}{Transposisi} membalik urutan hasil kali:
$(vu)^{\mathsf T} = u^{\mathsf T}v^{\mathsf T}$; melupakan
pembalikan itu merusak setiap perhitungan yang melibatkan invers. (iii)
\omterm{def:b2:linalg:annihilator}{Anihilator} $F^\circ$ hidup di $E^*$, bukan di $E$: ia menjadi
``komplemen ortogonal'' yang lazim itu hanya setelah suatu hasil kali
dalam menyamakan $E$ dengan $E^*$ (\cref{ch:b2:quadratic}); tidak ada
penyamaan semacam itu yang kanonik. (iv) ``Rank baris sama dengan rank
kolom'' tidak berarti \emph{ruang} baris sama dengan ruang kolom ---
keduanya hidup di ruang yang berbeda ($K^n$ dan $K^m$) dan
berhubungan lewat \cref{prop:b2:linalg:transposerank}, bukan sama.
(v) Rumus permutasi adalah alat bukti: untuk bilangan, pakailah
operasi baris dan kofaktor (\cref{ex:b2:linalg:sarrus}).
\end{remark}

\begin{example}[Pasangan trace membelah ruang matriks]\label{ex:b2:linalg:tracepairing}
Pada $\mathcal{M}_2(\R)$ dengan pasangan $\langle A, B\rangle =
\operatorname{tr}(AB)$ dari \cref{exo:b2:linalg:9}: uraikan $M =
\begin{psmallmatrix}1 & 4\\ 2 & 3\end{psmallmatrix}$ menjadi
bagian setangkup dan bagian antisetangkup,
\[
M = S + A, \qquad
S = \tfrac12(M + M^{\mathsf T}) =
\begin{pmatrix}1 & 3\\ 3 & 3\end{pmatrix},
\qquad
A = \tfrac12(M - M^{\mathsf T}) =
\begin{pmatrix}0 & 1\\ -1 & 0\end{pmatrix}.
\]
Maka $\operatorname{tr}(SA) = \operatorname{tr}
\begin{psmallmatrix}-3 & 1\\ -3 & 3\end{psmallmatrix} = 0$: kedua
bagiannya ``ortogonal'' terhadap pasangan trace --- sebuah
perwujudan fakta umum (yang dibuktikan pada soal akhir pekan
bab ini) bahwa matriks antisetangkup persis membentuk
\omterm{def:b2:linalg:annihilator}{anihilator} matriks setangkup. Dualitas melihat penguraian
$\mathcal{M}_n = \mathcal{S}_n \oplus
\mathcal{A}_n$ sebelum hasil kali dalam apa pun dipilih.
\end{example}

\begin{remark}[Pandangan ke depan di dalam jilid ini]\label{rem:b2:linalg:perspectives}
Perhatikan ketiga konstruksi bab ini berganti kostum di
halaman-halaman berikutnya. \emph{\omterm{def:b2:linalg:transpose}{Transpos}} kembali pada
\cref{ch:b2:reduction}: $u$ dan $u^{\mathsf T}$ berbagi
nilai eigen dengan multiplisitas geometrik yang sama (soal akhir
pekan bab ini, pertanyaan 15), dan itulah sebabnya telaah baris dan
telaah kolom sebuah matriks tak pernah berselisih. \emph{\omterm{def:b2:linalg:det}{Determinan}}
menjadi fungsi sebuah parameter pada \cref{ch:b2:reduction}
($\chi_u(X) = \det(X\,\mathrm{id} - u)$) dan menjadi Jacobian pada
\cref{ch:b2:multint}, tempat kemultilinearannya berubah menjadi
faktor penggantian peubah. \emph{Trace} menyemai
invarian keserupaan: ia koefisien kedua pada
$\chi_u$, jumlah nilai eigen, dan pada akhirnya integral
diagonal pada kesamaan bergaya \cref{ch:b2:fourier}. Satu
bab aljabar linear, tiga bayangan panjang.
\end{remark}

\begin{remark}[Di mana bab ini dipakai]
\omterm{def:b2:linalg:dual}{Ruang dual} bukan abstraksi demi abstraksi: \omterm{def:b2:linalg:annihilator}{anihilator}
dan \omterm{def:b2:linalg:transpose}{transpos} menjalankan teori keterselesaian sistem linear (soal
akhir pekan bab ini membuktikan alternatif Fredholm berdimensi hingga
dari keduanya), pasangan tak degenerat muncul kembali sebagai
bentuk polar pada \cref{ch:b2:quadratic} dan sebagai adjoin pada
\cref{ch:b2:hermitian}, dan \omterm{def:b2:linalg:det}{determinan} yang \omterm{def:b2:structures:generated}{dibangun} di sini
menggerakkan seluruh \cref{ch:b2:reduction}. Pada jilid Tahun ke-3
dualitas yang sama, setelah diangkut ke dimensi tak hingga, menjadi
teorema representasi Riesz dan teori Fredholm pada ruang
Hilbert --- dengan kekompakan menggantikan pencacahan dimensi yang
dipakai di sini.
\end{remark}

\section{Latihan}

\begin{exercise}[$\star$]\label{exo:b2:linalg:1}
Di $\R^3$, misalkan $\varphi_1(x,y,z) = x + y$, $\varphi_2 = y + z$,
$\varphi_3 = x + z$. Buktikan bahwa $(\varphi_1, \varphi_2, \varphi_3)$
adalah basis $(\R^3)^*$ lalu carilah basis $\R^3$ yang dualnya adalah
basis itu.
\end{exercise}

\begin{exercise}[$\star$]\label{exo:b2:linalg:2}
Hitung \omterm{def:b2:linalg:det}{determinan} berikut dengan rumus permutasi:
\[
\begin{pmatrix} 0 & 0 & a\\ 0 & b & 0\\ c & 0 & 0 \end{pmatrix},
\qquad
\begin{pmatrix}
a & b & 0 & 0\\
c & d & 0 & 0\\
0 & 0 & e & f\\
0 & 0 & g & h
\end{pmatrix},
\]
lalu nyatakan aturan diagonal blok yang disarankan matriks kedua.
\end{exercise}

\begin{exercise}[$\star$]\label{exo:b2:linalg:3}
Misalkan $F = \{(x,y,z,t) \in \R^4 : x + y = z + t \text{ dan } x = 2y\}$.
Berilah sebuah basis $F^\circ$ lalu periksa
\cref{thm:b2:linalg:annihilator} pada dimensinya.
\end{exercise}

\begin{exercise}[$\star\star$]\label{exo:b2:linalg:4}
Misalkan $a_0, \dots, a_n$ titik yang berbeda di $K$ dan $\varphi_i \colon
P \mapsto P(a_i)$ pada $K_n[X]$. Buktikan bahwa $(\varphi_0, \dots,
\varphi_n)$ adalah basis $K_n[X]^*$, kenali basis pra-dualnya,
lalu uraikan bentuk $P \mapsto \int_0^1 P(t)\,\dd t$ (untuk $K = \R$,
$n = 2$, $a_i = 0, \frac12, 1$) pada basis itu --- sehingga aturan
Simpson dikenali.
\end{exercise}

\begin{exercise}[$\star\star$]\label{exo:b2:linalg:5}
Misalkan $u \in \mathcal{L}(E)$ dengan $\dim E = n$ dan $\operatorname{rk} u
= 1$. Buktikan bahwa $u = \varphi(\cdot)\, a$ untuk suatu vektor $a$ dan
suatu bentuk $\varphi$; bahwa $\operatorname{tr} u = \varphi(a)$; dan bahwa
$u^2 = (\operatorname{tr} u)\, u$. Turunkan $\det(I + u) = 1 +
\operatorname{tr} u$.
\end{exercise}

\begin{exercise}[$\star\star$]\label{exo:b2:linalg:6}
Buktikan bahwa setiap hiperbidang $\mathcal{M}_n(K)$ ($n \geq 2$)
memuat sebuah matriks yang punya invers. \emph{Petunjuk: sebuah
hiperbidang berbentuk $\{M : \operatorname{tr}(AM) = 0\}$ untuk suatu
$A \neq 0$ (\cref{exo:b2:linalg:9}). Jika $A$ skalar, tunjukkan sebuah
matriks bertrace nol yang punya invers; jika tidak, carilah $M$ berinvers
yang membuat diagonal $AM$ nol --- matriks bergaya permutasi sanggup.}
\end{exercise}

\begin{exercise}[$\star\star$]\label{exo:b2:linalg:7}
(Turunan \omterm{def:b2:linalg:det}{determinan}) Untuk $A \in \mathcal{M}_n(\R)$, buktikan
dari kemultilinearan bahwa
\[
\frac{\dd}{\dd t}\Big|_{t=0} \det(I_n + tA) = \operatorname{tr} A ,
\]
lalu turunkan $\det(\eu^{tA}) = \eu^{t\operatorname{tr} A}$ dengan
mengandaikan keterdiferensialan $t \mapsto \det(\eu^{tA})$ dan sifat
grup $\eu^{(s+t)A} = \eu^{sA}\eu^{tA}$ (yang ditegakkan pada
\cref{ch:b2:diffeq}).
\end{exercise}

\begin{exercise}[$\star\star$]\label{exo:b2:linalg:8}
(Sirkulan $3 \times 3$) Misalkan $j = \eu^{2\iu\pi/3}$ dan
\[
C = \begin{pmatrix}
a & b & c\\
c & a & b\\
b & c & a
\end{pmatrix} \in \mathcal{M}_3(\C).
\]
Periksalah bahwa kolom matriks Vandermonde atas $1, j, j^2$ adalah
vektor eigen $C$, lalu turunkan
\[
\det C = (a + b + c)(a + bj + cj^2)(a + bj^2 + cj).
\]
\end{exercise}

\begin{exercise}[$\star\star\star$]\label{exo:b2:linalg:9}
Buktikan bahwa setiap bentuk linear $t$ pada $\mathcal{M}_n(K)$ berbentuk
$M \mapsto \operatorname{tr}(AM)$ untuk suatu $A$ yang tunggal: pemetaan
$A \mapsto \operatorname{tr}(A\,\cdot)$ adalah isomorfisma dari
$\mathcal{M}_n(K)$ pada dualnya. Turunkan lagi pernyataan ketunggalan
pada \cref{prop:b2:linalg:trace}.
\end{exercise}

\begin{exercise}[$\star\star\star$]\label{exo:b2:linalg:10}
Misalkan $u, v \in \mathcal{L}(E)$ dengan $u \circ v - v \circ u = u$.
Buktikan bahwa $u$ nilpoten. \emph{Petunjuk: tunjukkan
$\operatorname{tr}(u^k) = 0$ untuk setiap $k \geq 1$ (hitung $u^k v -
v u^k$ secara induktif), lalu pakai fakta berikut, yang dibuktikan lewat
kesamaan Newton atau lewat induksi pada dimensinya: endomorfisma
sebuah ruang vektor atas $\C$ yang semua pangkatnya bertrace nol
pastilah nilpoten. Bekerjalah atas $\C$.}
\end{exercise}

\begin{exercise}[$\star\star$]\label{exo:b2:linalg:11}
(Vandermonde) Untuk $a_0, \dots, a_n \in K$, buktikan
\[
\det\begin{pmatrix}
1 & 1 & \cdots & 1\\
a_0 & a_1 & \cdots & a_n\\
\vdots & \vdots & & \vdots\\
a_0^n & a_1^n & \cdots & a_n^n
\end{pmatrix}
= \prod_{0 \leq i < j \leq n} (a_j - a_i).
\]
\emph{(Pandang \omterm{def:b2:linalg:det}{determinan} itu sebagai polinomial dalam $a_n$: kenali
derajatnya, akarnya, dan koefisien utamanya; lalu berinduksilah.)}
\end{exercise}

\begin{exercise}[$\star\star\star$]\label{exo:b2:linalg:12}
Misalkan $A, B, C, D \in \mathcal{M}_n(K)$ dengan $K$ tak hingga, dan
andaikan $CD = DC$. Buktikan bahwa
\[
\det\begin{pmatrix} A & B\\ C & D\end{pmatrix}
= \det(AD - BC).
\]
\emph{(Tinjau dulu $D$ yang punya invers, dengan mengalikan dari kanan
oleh $\begin{psmallmatrix} I & 0\\ -D^{-1}C & I\end{psmallmatrix}$;
lalu ganti $D$ dengan $D + tI$ dan bandingkan dua polinomial dalam $t$.)}
\end{exercise}

\section{Soal: Alternatif Fredholm}

Kapan sistem linear $u(x) = b$ punya penyelesaian? Jawaban
lengkapnya berupa pernyataan dualitas: \emph{tepat ketika $b$
dianihilasi oleh setiap bentuk linear yang menganihilasi peta
$u$} --- dan bentuk semacam itu terhitungkan, sebab ia kernel
transposnya. Soal akhir pekan ini membangun kamus lengkap
dualitas berdimensi hingga (pemfaktoran bentuk, bidualitas,
kalkulus \omterm{def:b2:linalg:annihilator}{anihilator}, \omterm{def:b2:linalg:transpose}{transpos}), membuktikan \emph{alternatif
Fredholm} berdimensi hingga, dan ditutup dengan
bentuk trace beserta sebuah pencirian: trace adalah satu-satunya
invarian linear bagi keserupaan. Di sepanjang soal ini, $E$ dan $F$
adalah ruang vektor atas $K$ yang berdimensi hingga, dan $n = \dim E$.

\begin{problem}[{Soal akhir pekan --- dualitas dalam dimensi hingga
dan alternatif Fredholm}]
\label{pb:b2:linalg:1}
Notasi: untuk $S \subseteq E^*$, \emph{\omterm{def:b2:linalg:annihilator}{pra-anihilator}} adalah
$S_\circ = \{x \in E : \varphi(x) = 0 \text{ untuk setiap } \varphi
\in S\}$; adapun \omterm{def:b2:linalg:annihilator}{anihilator} $F^\circ$ dan \omterm{def:b2:linalg:transpose}{transpos} $u^{\mathsf T}$
adalah yang ada pada \cref{def:b2:linalg:annihilator} dan
\cref{def:b2:linalg:transpose}.

\medskip
\noindent\textbf{Bagian I --- Lema pemfaktoran.} Misalkan
$\varphi_1, \dots, \varphi_p, \varphi \in E^*$.
\begin{enumerate}
  \item Misalkan $\Phi \colon E \to K^p$, $x \mapsto (\varphi_1(x),
        \dots, \varphi_p(x))$. Kenali $\ker\Phi$, tunjukkan
        $\Phi^{\mathsf T}$ memetakan bentuk koordinat $K^p$ ke
        $\varphi_i$, lalu turunkan
        \[
        \dim \bigl(\ker\varphi_1 \cap \dots \cap
        \ker\varphi_p\bigr) = n - \dim
        \operatorname{Vect}(\varphi_1, \dots, \varphi_p).
        \]
  \item (Lema pemfaktoran) Buktikan kesetaraan berikut:
        \[
        \varphi \in \operatorname{Vect}(\varphi_1, \dots,
        \varphi_p)
        \iff
        \ker\varphi_1 \cap \dots \cap \ker\varphi_p \subseteq
        \ker\varphi .
        \]
  \item Turunkan: $(\varphi_1, \dots, \varphi_p)$ bebas bila dan hanya
        bila $\bigcap_i \ker\varphi_i$ berdimensi $n - p$; dan sebuah
        subruang berkodimensi $p$ adalah irisan $p$
        hiperbidang, tak pernah kurang.
  \item Di $\R^4$, misalkan $\varphi_1 = x + y - z$, $\varphi_2 = y +
        z - t$, $\psi = x + 2y - t$ dan $\psi' = x + y + t$.
        Putuskan, dengan lema pemfaktoran, apakah $\psi$ dan
        $\psi'$ termasuk $\operatorname{Vect}(\varphi_1,
        \varphi_2)$.
  \item Pada $E = \R_2[X]$, tunjukkan bahwa $\psi_0 \colon P \mapsto
        P(0)$, $\psi_1 \colon P \mapsto P(1)$, $\psi_2 \colon P
        \mapsto \int_0^1 P(t)\dd t$ membentuk basis $E^*$,
        hitunglah basis $(P_0, P_1, P_2)$ di $E$ yang dualnya adalah
        basis itu, lalu carilah satu-satunya $P \in \R_2[X]$ dengan $P(0)
        = 1$, $P(1) = 2$, $\int_0^1 P = \frac32$.
\end{enumerate}

\medskip
\noindent\textbf{Bagian II --- Bidualitas dan kalkulus
\omterm{def:b2:linalg:annihilator}{anihilator}.}
\begin{enumerate}[resume]
  \item Tunjukkan bahwa \emph{pemetaan evaluasi} $J \colon E \to
        E^{**}$, $J(x)(\varphi) = \varphi(x)$, bersifat linear dan
        injektif, jadi merupakan isomorfisma dalam dimensi hingga.
  \item (\omterm{def:b2:linalg:annihilator}{Anihilator} ganda) Tunjukkan $J(F) = F^{\circ\circ} :=
        (F^\circ)^\circ$ untuk setiap subruang $F \subseteq E$: di
        bawah penyamaan $J$, \omterm{def:b2:linalg:annihilator}{anihilator} dari \omterm{def:b2:linalg:annihilator}{anihilator}
        adalah subruangnya sendiri.
  \item Buktikan kalkulus \omterm{def:b2:linalg:annihilator}{anihilator}: $(F + G)^\circ = F^\circ
        \cap G^\circ$ dan $(F \cap G)^\circ = F^\circ + G^\circ$.
  \item Turunkan (lalu buktikan ulang secara langsung): dua bentuk tak
        nol yang kernelnya sama pastilah sebanding.
  \item (Basis antedual) Tunjukkan bahwa untuk setiap basis $(\varphi_1,
        \dots, \varphi_n)$ di $E^*$ ada tepat satu basis
        $(u_1, \dots, u_n)$ di $E$ dengan $\varphi_i(u_j) =
        \delta_{ij}$.
\end{enumerate}

\medskip
\noindent\textbf{Bagian III --- Kalkulus \omterm{def:b2:linalg:transpose}{transpos}.}
\begin{enumerate}[resume]
  \item Tunjukkan bahwa $u \mapsto u^{\mathsf T}$ adalah bijeksi linear
        dari $\mathcal{L}(E, F)$ pada $\mathcal{L}(F^*, E^*)$,
        dan bahwa $(u^{-1})^{\mathsf T} = (u^{\mathsf T})^{-1}$
        bila $u$ punya invers.
  \item (Kealamian) Tunjukkan bahwa $u^{\mathsf T\mathsf T} \circ J_E
        = J_F \circ u$: di bawah isomorfisma evaluasi,
        \omterm{def:b2:linalg:transpose}{transpos} ganda \emph{adalah} $u$.
  \item Tunjukkan: $u$ surjektif bila dan hanya bila $u^{\mathsf T}$
        injektif; dan $u$ injektif bila dan hanya bila $u^{\mathsf T}$
        surjektif.
  \item Untuk $u \in \mathcal{L}(E)$: sebuah subruang $F$ stabil
        terhadap $u$ bila dan hanya bila $F^\circ$ stabil terhadap
        $u^{\mathsf T}$.
  \item Tunjukkan bahwa $\ker(u^{\mathsf T} - \lambda\,
        \mathrm{id}_{E^*}) = \bigl(\operatorname{im}(u - \lambda\,
        \mathrm{id}_E)\bigr)^\circ$, lalu turunkan bahwa $u$ dan
        $u^{\mathsf T}$ punya nilai eigen yang sama dengan
        multiplisitas geometrik yang sama.
\end{enumerate}

\medskip
\noindent\textbf{Bagian IV --- Alternatif Fredholm.}
\begin{enumerate}[resume]
  \item Buktikan bahwa $\operatorname{im} u = (\ker u^{\mathsf
        T})_\circ$ untuk $u \in \mathcal{L}(E, F)$, lalu turunkan
        \emph{alternatif Fredholm} dalam dimensi hingga: persamaan
        $u(x) = b$ punya penyelesaian bila dan hanya bila setiap
        $\psi \in F^*$ dengan $u^{\mathsf T}\psi = 0$ memenuhi
        $\psi(b) = 0$.
  \item Bentuk matriksnya: untuk $A \in \mathcal{M}_{m,n}(K)$ dan $b \in
        K^m$, tepat satu dari berikut ini berlaku: (i) $Ax = b$
        punya penyelesaian; (ii) ada $y \in K^m$ dengan
        $A^{\mathsf T}y = 0$ dan $y^{\mathsf T}b = 1$. Buktikan baik
        ``paling banyak satu'' maupun ``sedikitnya satu''.
  \item Carilah semua $b \in \R^3$ yang membuat sistem
        \[
        x + y = b_1, \qquad y + z = b_2, \qquad x + 2y + z = b_3
        \]
        punya penyelesaian, dengan menghitung kernel matriks
        transposnya.
  \item (Sebuah masalah Neumann diskret) Pada $E = \R^n$ ($n \geq 3$),
        tetapkan $L$ oleh $(Lx)_k = x_k - \frac12(x_{k-1} +
        x_{k+1})$, dengan indeks modulo $n$. Tunjukkan $L^{\mathsf T} = L$
        (dengan penyamaan kanonik), tunjukkan $\ker L$ adalah garis
        vektor konstan \emph{(perhatikan sebuah koordinat maksimal)},
        lalu simpulkan: $Lx = b$ terselesaikan bila dan hanya bila
        $\sum_k b_k = 0$.
\end{enumerate}

\medskip
\noindent\textbf{Bagian V --- Bentuk trace dan teorema
keawetan.} Ingat kembali dari \cref{exo:b2:linalg:9} bahwa $A \mapsto
\operatorname{tr}(A\,\cdot)$ menyamakan $\mathcal{M}_n(K)$ dengan
dualnya. Andaikan $\operatorname{char} K = 0$ (misalnya $K = \Q, \R,
\C$).
\begin{enumerate}[resume]
  \item Di bawah penyamaan itu, tunjukkan bahwa \omterm{def:b2:linalg:annihilator}{anihilator}
        subruang $\mathcal{S}_n$ berisi matriks setangkup adalah
        subruang $\mathcal{A}_n$ berisi matriks antisetangkup, dan
        sebaliknya.
  \item Tunjukkan bahwa \omterm{def:b2:linalg:annihilator}{anihilator} hiperbidang
        $\mathfrak{sl}_n = \{M : \operatorname{tr} M = 0\}$ adalah
        garis $K I_n$; setara dengan itu, bentuk linear yang nol
        pada semua matriks bertrace nol pastilah kelipatan tracenya.
  \item Tunjukkan bahwa setiap matriks $\mathcal{M}_n(K)$ adalah jumlah
        dua matriks yang punya invers.
  \item (Trace satu-satunya invarian linear bagi keserupaan) Misalkan
        $t$ sebuah bentuk linear pada $\mathcal{M}_n(K)$ dengan
        $t(PMP^{-1}) = t(M)$ untuk setiap $M$ dan setiap $P$ yang
        punya invers. Tunjukkan lebih dulu $t(PX) = t(XP)$ untuk $P$
        berinvers, lalu $t(BX) = t(XB)$ untuk \emph{setiap} $B$, dan
        simpulkan $t = c \operatorname{tr}$ untuk suatu $c \in K$.
  \item Tunjukkan bahwa $\operatorname{rk} u \leq r$ bila dan hanya bila
        $u$ adalah jumlah $r$ pemetaan berank $\leq 1$, yakni $u =
        \sum_{i=1}^{r} \psi_i(\cdot)\,f_i$ dengan $\psi_i \in E^*$,
        $f_i \in F$; lalu turunkan $\operatorname{rk}(u + v) \leq
        \operatorname{rk} u + \operatorname{rk} v$.
  \item (Rangkuman) Susunlah kamus yang dibuktikan pada soal
        ini: subruang lawan \omterm{def:b2:linalg:annihilator}{anihilator}, jumlah lawan
        irisan, pemetaan lawan \omterm{def:b2:linalg:transpose}{transpos}, keterselesaian lawan
        keortogonalan terhadap kernel \omterm{def:b2:linalg:transpose}{transpos}, trace lawan
        keserupaan. Untuk tiap entri, sebutkan pertanyaan yang
        membuktikannya, dan nyatakan dalam satu kalimat apa yang
        menggantikan pencacahan dimensi ketika dimensinya menjadi tak
        hingga (jilid Tahun ke-3 memerincinya pada ruang Hilbert).

\end{enumerate}
\end{problem}
